\section{模型假设与符号说明}
\label{sec:assume}

\subsection{模型假设}

\begin{enumerate}[label=\textbf{假设\arabic*.}, leftmargin=4.2em, itemsep=2pt]
\item 各时段附件给出的功率（kW）为该 10 分钟区间的平均功率，故区间电量（kWh）$=$ 功率 $\times \frac{10}{60}$；
\item 光伏发电优先本地消纳，富余功率可用于储能充电；当富余功率超过储能最大充电功率时允许\textbf{弃光}；
\item 储能充放电效率取\textbf{单程 90\%}（充电存入 90\%、放电放出 90\%，往返 81\%），且同一时段不会同时充放电（LP 最优解自动满足）；
\item 微网与外网的联络线购电能力不受限（无功率上限），但供电不得低于小区负载，缺口须以 5 倍电价紧急购电；
\item 问题二中 0:00 制定计划时\textbf{无法获得当日实测负载与光伏}，只能利用历史实测数据构造预测；
计划量报出后不可更改，日内由储能按第 \ref{sec:q2} 节的\textbf{实时平衡控制规则}兜底
（只用当前实测与当前储电量），兜不住的缺口才按 5 倍电价紧急购电；
\item 问题三中 0:00 计划依据 0:00 发布的光伏预报，6:00/12:00/18:00 的调整仅作用于其后时段并继承当前储电量；
\item 问题四中规划时已知当日实时电价（现货日前市场口径）；
\item 结果文件窗口为 2025-02-01 至 12-31；2025 年 1 月作为储电量链条与预测历史信息的预热期（初始储电量 6000 kWh 对应 2025-01-01 0:00）；
\item 日末储电量以“次日平均电价 $\times$ 放电效率”作为残值计入当日目标函数（滚动时域的标准做法，避免每日耗尽储能的短视行为）；
\item 忽略线路损耗、储能自放电与设备老化，各时段参数在日内不变。
\end{enumerate}

\subsection{符号说明}

\begin{table}[htbp]
\centering
\caption{主要符号}
\label{tab:symbols}
\small
\begin{tabular}{@{}clcl@{}}
\toprule
符号 & 含义 & 符号 & 含义 \\
\midrule
$t$ & 时段序号，$t=1,\dots,144$ & $C_t,D_t$ & 时段 $t$ 的充电量、放电量（kWh）\\
$L_t$ & 时段 $t$ 小区负载电量（kWh） & $E_t$ & 时段 $t$ 起始储电量（kWh）\\
$PV_t$ & 时段 $t$ 光伏发电量（kWh） & $E_0,E_{145}$ & 0:00、24:00 储电量（kWh）\\
$P_t$ & 时段 $t$ 计划购电量（kWh） & $\eta_c,\eta_d$ & 充、放电效率，均取 0.9\\
$p_t$ & 时段 $t$ 电价（元/kWh） & $v$ & 日末储电量残值系数（元/kWh）\\
$q$ & 报童安全裕度的分位水平 & $b_t$ & 时段 $t$ 的安全裕度（kWh）\\
$\hat L_t,\hat{PV}_t$ & 负载、光伏的预测（基值） & $S_t,X_t$ & 计划超出量、调整超出量（kWh）\\
$\bar C,\bar D$ & 充放电功率上限折合 833.33 kWh & $E_t^{\text{急}}$ & 时段 $t$ 紧急购电量（kWh）\\
\bottomrule
\end{tabular}
\end{table}

\subsection{通用单日线性规划模型}

综合第 2 节的机理与第 3 节的规律，四问共用的单日 LP 如下（以预测基值 $\hat L,\hat{PV}$ 为输入）：

\begin{equation}
\underset{P,C,D,E}{\text{min}}\quad \sum_{t=1}^{144} p_t P_t \;-\; v\,E_{145}
\label{eq:obj}
\end{equation}
\begin{align}
\text{s.t.}\quad & \hat{PV}_t + P_t + D_t - C_t \ \ge\ \hat L_t, && t=1,\dots,144 \label{eq:c1}\\
& E_{t+1} = E_t + \eta_c C_t - \frac{D_t}{\eta_d}, && t=1,\dots,144 \label{eq:c2}\\
& 1200 \le E_t \le 10800, && t=1,\dots,145 \label{eq:c3}\\
& 0 \le C_t \le \bar C,\quad 0 \le D_t \le \bar D, && t=1,\dots,144 \label{eq:c4}\\
& P_t \ge 0, && t=1,\dots,144 \label{eq:c5}
\end{align}
其中 $E_1=E_0$ 由链条给出（问题一取 $E_1=E_{145}=6000$；问题二至四取前一日的 24:00 储电量），
残值系数 $v=0.9\times\bar p_{\text{次日}}$（问题一不含残值项而直接约束端点相等）。
该模型为\textbf{标准线性规划}，决策变量 $4\times144+1=577$ 个，约束 $288+1$ 条，
用 HiGHS 内点/单纯形混合算法求解，单日求解时间约 9 ms。
